% Generated by scripts/assemble.py; corrections belong in source/corrections.json.
\examhead{115}{厦门大学 2026 年数学分析试卷}{112}
% p113-b08; source PDF page 113
\blockanchor{p113-b08}
\question{1}(15 分) 设  $a_{n} > 0 (n = 1, 2, \cdots)$ ，级数  $\sum_{n=1}^{\infty} a_{n}$  发散，证明： $\sum_{n=1}^{\infty} \frac{a_{n}}{S_{n}}$  发散，其中  $S_{n} = a_{1} + a_{2} + \cdots + a_{n}$ .\par
% p113-b09; source PDF page 113
\blockanchor{p113-b09}
\question{2}(15 分) 设  $x_{n}=1+\frac{1}{2}+\cdots+\frac{1}{n}-\ln(n+1)(n=1,2,\cdots)$ ，证明： $\lim_{n\to\infty}x_{n}$  存在.\par
% p113-b10; source PDF page 113
\par\Needspace{6\baselineskip}
\blockanchor{p113-b10}
\question{3}(20 分) 设 $f(x)$ 是定义在 $[0,1]$ 上的连续可微函数, 且满足 $f(0) = 0$, $f(1) = 1$, 证明:\par
\nopagebreak[4]
% p113-b11; source PDF page 113
\blockanchor{p113-b11}
\displayformula{\int_ {0} ^ {1} | f (x) - f ^ {\prime} (x) | \mathrm{d} x \geq \frac {1}{e}.}
% p113-b12; source PDF page 113
\blockanchor{p113-b12}
\question{4}（20分）设 $f(x)$ 在 $(- \infty, + \infty)$ 上有三阶连续导数，且 $f(x), f'''(x)$ 在 $(- \infty, + \infty)$ 上有界，证明：$f'(x), f''(x)$ 在 $(- \infty, + \infty)$ 上有界.\par
% p113-b13; source PDF page 113
\par\Needspace{6\baselineskip}
\blockanchor{p113-b13}
\question{5}(20 分) 设 $f(x), g(x)$ 是定义在 $[0, +\infty)$ 上的非负连续函数, 满足\par
\nopagebreak[4]
% p113-b14; source PDF page 113
\blockanchor{p113-b14}
\displayformula{f (x) \leq \int_ {0} ^ {x} f (s) g (s) \mathrm{d} s, \forall x \geq 0.}
% p113-b15; source PDF page 113
\blockanchor{p113-b15}
证明： $f(x) \equiv 0, x \in [0, +\infty)$ .\par
% p113-b16; source PDF page 113
\par\Needspace{6\baselineskip}
\blockanchor{p113-b16}
\question{6}(20 分) 记 $\mathbb{R}^3$ 中的点 $x = (x_1, x_2, x_3)$, $\mathbb{R}^3$ 中以原点为圆心, $r > 0$ 为半径的开球记为 $B_r$, 其边界为 $\partial B_r$. 设 $f: \mathbb{R}^3 \to \mathbb{R}$ 具有连续的偏导数, 且 $f(x) = 0$, $\forall x \in \partial B_r$, 证明:\par
\nopagebreak[4]
% p113-b17; source PDF page 113
\blockanchor{p113-b17}
\displayformula{\lim _ {\varepsilon \rightarrow 0 ^ {+}} \iiint_ {B _ {r} \backslash B _ {\varepsilon}} \frac {\sum_ {i = 1} ^ {3} x _ {i} \frac {\partial f}{\partial x _ {i}} (\boldsymbol {x})}{| \boldsymbol {x} | ^ {3}} \mathrm{d} V = - 4 \pi f (\boldsymbol {0}).}
% p113-b18; source PDF page 113
\blockanchor{p113-b18}
\question{7}(20 分) 设 $V$ 是由球面 $x^{2} + y^{2} + z^{2} = 2az$ 与 $x^{2} + y^{2} + z^{2} = az$ 所围立体 $(a > 0)$, 计算 $\iiint_{V} z \, \mathrm{d}x \, \mathrm{d}y \, \mathrm{d}z$.\par
% p113-b19; source PDF page 113
\blockanchor{p113-b19}
\question{8}(20 分) 设  $f(x)$  在  $[-\pi , \pi ]$  上连续,  $a_{n} (n = 0, 1, \cdots)$ ,  $b_{n} (n = 1, 2, \cdots)$  为 f 的 Fourier 系数. 证明: 级数  $\sum_{n=1}^{\infty} \frac{a_{n}}{n}$  和  $\sum_{n=1}^{\infty} \frac{b_{n}}{n}$  均收敛.\par
